\section{Limitations}\label{sec:limitations}

\subsection{Internal validity}

The theory holds conditionally on the scalars and on the design. The study measures the learning of the scalars in two ways,
the linearized model and the full sampler, and the second only on three cells with 50 replicates each, because one full
posterior costs about \cbCpuHours~CPU-hours at the product's size. The closed forms use the plant's noise scale, which the full
method learns; the verification cells measure what that costs. The linearization is taken at the prior centre, which is the
shipped engine, and the virtual plants put the truth at the posterior mean of the ten surveys, a distance of up to one prior
standard deviation on the power scalars; a plant further from the prior centre would make the linearized predictor worse and
the verification would have to be repeated there. The orthogonal prior is implemented in its empirical form over the training
surveys, which is the finite-dimensional form of the orthogonality of Plumlee~\cite{plumlee} with respect to the empirical
distribution of the operating points; a version orthogonal with respect to the operating window as a whole would differ by
the difference between the two distributions. The scaled prior is used with one value of its scaling constant. The Student-$t$
bands of the regressions use the residual variance and are exact under Gaussian noise; the Gaussian bands of the closed forms
use the true noise scale, which favours them.

\subsection{External validity}

The virtual plants are models. V0 and V1 are built from the engine itself, so their curvature is the engine's curvature,
scaled; V2 uses the capacity law that Rocha et al.~\cite{rocha} fitted to the same ten surveys, which is curved in the setting
with the exponent 1.27 their Table~4 prints, and which is itself a fit. The study therefore says when the method can win,
and by how much, on plants that curve as these do; it does not say that a given industrial crusher curves this way. A
three-point series of an HP400 in the work of Duarte et al.~\cite{duarte} shows a capacity concave in the setting over 6~mm,
which is evidence that real capacities can curve, not a measurement of how much. No discrete-element simulation is used as a
truth: one operating point of a quarter section of a secondary cone crusher took 172 to 404 CPU-hours~\cite{quistthesis}, and
the study needs thousands. The throw and the speed of the HP500 are held at 34~mm and 360~rpm, product assumptions for which no
primary source was found; the design with varied speed moves it by ten percent either way around the engine's capacity
optimum, and its conclusion depends on that range.

\subsection{Construct validity}

The risk is the mean squared error of a point prediction, and the comparison with the empirical predictors is made on it and
on the mean absolute percentage error; a decision that weighs an overprediction of power differently from an underprediction
would need another loss, under which Theorem~\ref{thm:risk} still gives the error's law under Gaussian noise but not the
ranking. The band is judged by its coverage and its width at the 80\pct{} level only. The discrepancy is a plane in two
regressors, as in the released method; a Gaussian-process discrepancy would change Proposition~\ref{prop:bound} (its kernel
would absorb curved sensitivities too) and the results of Section~\ref{sec:results-capacity}. The plants assume one ore and one
machine state across the surveys; a drift of the ore's hardness or of the liners' wear between surveys is a regressor that
neither the plant nor the predictors carry, and in a real set of surveys it would enter the noise and lower every gain measured here.

\section{Reproducibility}\label{sec:repro}

Every number, table and figure of this paper is computed from the files of the repository~\cite{chancadem} at release
\productrelease{} by committed scripts. The harness is \path{manuscripts/crusherbayes-theory/sim/}: \path{run_exact.py} runs
the closed-form study, \path{run_design.py} the design study and \path{run_sampler.py} the simulation-based calibration and the
verification; the identifiability analysis and the pilot are in \path{manuscripts/crusherbayes-theory/analysis/}; the
figures, the tables and the numbers of the text are written by \path{manuscripts/crusherbayes-theory/figures/make_figs.py}
from the summaries in \path{manuscripts/crusherbayes-theory/data/}. The tests of the harness, run with \texttt{pytest},
check the closed forms against the released implementation of the method, every identity of Section~\ref{sec:theory} on random
instances, and the independence of the results from the number of processes. The engine and the method are those of the
release, imported from \path{data-pipeline/pipeline/}; nothing in them was changed for this paper. The ten surveys are the
open table of Rocha et al.~\cite{rocha} (CC~BY~4.0), committed as \path{data/examples/hp500-surveys.csv}. The run times on the
machine of the study (\envCpus{} logical processors, \envPlatform, Python~\envPython, numpy~\envNumpy, emcee~\envEmcee), recorded
by \path{analysis/record_runs.py} in \path{data/run_times.json}, were \runExactTime{} for the closed-form study,
\runDesignTime{} for the design study, \runIdentTime{} for the identifiability analysis, \runTauTime{} for the analysis of the
discrepancy scale and the coverage, and \runSamplerTime{} for the full-sampler runs.

\section{Conclusions}\label{sec:conclusions}

Given its scalars, the prediction of a Bayesian calibration with a discrepancy linear in the regressors is least squares on
those regressors, shrunk toward the engine's plane, plus the engine's departure from that plane, and its risk has the exact form
of Theorem~\ref{thm:risk}. Over an empirical regression in the same regressors the engine gains in two ways: a shrinkage gain,
bounded by the variance of the fitted plane and fading with the number of surveys, and a curvature gain, which persists and is
positive only when the engine's curvature is aligned with the plant's and less than twice as large. On \cellsTotal{} cells of
virtual plants the theory matched the simulation to Monte Carlo error. For the HP500 engine of ChancaDEM the capacity is affine in
the setting, so on throughput only the shrinkage acts; on power the engine's curvature is \curvShareKw\pct{} of the noise
variance at best. A curved capacity of the plant's form gains a persistent \capRoRmseS{} to \capRoRmseW\pct{} of RMSE inside the
operating windows (\capRoMseS{} to \capRoMseW\pct{} of the mean squared error) and \capRoRmseX\pct{} under extrapolation, bounded
inside the training window by the plant's own curvature. The capacity base cannot be narrowed below a floor by more surveys at
the same speed over the same setting range, and a product-size measurement is the design that teaches the scalars most. The next steps are a product-size
measurement on the HP500 surveys, a capacity law curved in the setting inside the engine, and a discrepancy whose prior scale is
set by~\eqref{eq:taustar} rather than by convention.
